% Pursuing Stacks (cotes 134-2 to 134-8): archive cote 134-4, site page 95 (PDF page 96), typescript page 182.
% typescript. Draft compared: ps3.tex:4864-4921.
% AI-transcribed by Claude Opus 5.5 (claude-opus-5-5, adaptive thinking, effort high), Message Batches API, batch ps2 (msgbatch_01NVwr1Yuo8qJf5rykeGBXPK), 2026-11-12; not yet checked by eye.
\scanpage{134-4/95}
\mspage{182}

for the map f: $X\longrightarrow Y$ (at least any base change within the given context,
namely either spaces or schemes with etale topology), the corresponding
map \xout{f' should be aspheric. When this condition is satisfied, we}
\[ f' : \quad X'=X\times_{Y}Y'\longrightarrow Y' \]
is a weak equivalence, or what amounts to the same, that for any $Y'$ the
corresponding map is aspheric. This property, \xout{should} if a name is needed
indeed, would properly be called ``\uline{universally aspheric}'' . Thus in (Cat),
a map $f: X\to Y$ will be called either aspheric,\unsure{of|or} universally aspheric,
\deleted{if}\inserted{when} for any base change of the special type $\over{X}{Y}_{/\over{x}{y}}\to\over{X}{Y}$ namely ``localization''
in the first case, or any base change whatever $Y'\to Y$ in the second
case, the corresponding map $f'$ is a weak equivalence. On the other
hand, if $Y$ is just the final object of (Cat), it turns out that the
two notions for $X$ (being ``aspheric over $e$'' and being ``universally as-%
pheric over $e$'' ) coincide, and just mean that $X\to e$ is a weak equiva-%
lence. In accordance with the use which has been prevalent for a long
time\xout{,} in the context of spaces,\xout{(and which I} such an object will be
called simply an \uline{aspheric object} - which means that the corresponding
topos is aspheric (namely has ``trivial'' cohomology invariants, and
hence trivial homotopy invariants of any kind ...).

In case the notion of weak equivalence is replaced by a basic locali-%
zer\xout{,} $\underline{W}\subset F\unsure{\ell}((Cat))$, there is no reason whatever to change anything in
this terminology - \inserted{except that,} if needed, we will add the qualifying $\underline{W}$ , and
speak of \uline{$\underline{W}$-aspheric} or \uline{universally} \uline{$\underline{W}$-aspheric maps} in (Cat), as well as
of \uline{$\underline{W}$-aspheric} \uline{objects} of (Cat).

What about terminology for maps and objects within a category $A^{\wedge}$\hand{?}
Here the emphasis should be of course upon perfect coherence \xout{between}\inserted{with}
the \over{p}{t}erminology used in (Cat). An object F of $A^{\wedge}$ should always be sensed
in terms of the induced topos $A^{\wedge}_{/F}\hand{$\simeq$}(A_{/F})^{\wedge}$, or what amounts to the same\hand{,}
in terms of the corresponding object $A_{/F}$ in (Cat), which will imply \xout{th}
that `` F \inserted{is} \uline{aspheric} '' cannot possibly mean anything else but $A_{/F}$ \inserted{is} aspheric
as an object of (Cat) ; the same if qualifying by a $\underline{W}$ - F is called
$\underline{W}$ - \uline{aspheric} if $A_{/F}$ is a $\underline{W}$-aspheric object of (Cat). Similarly\xout{,} for
maps - thus $f: F\to G$ will be called a weak equivalence, if the
corresponding map for the induced topoi is a weak equivalence, or equi-%
valently, if the corresponding \uline{map in} (Cat)
\[ A_{/F}\longrightarrow \deleted{\illegible{1}}\ A_{/G} \]
is a weak equivalence. When a $\underline{W}$ is given, we would say instead (if con-%
fusion may arise) that f is a \uline{$\underline{W}$-equivalence}. The map will be called

\margin{95}
